\begin{abstract}
Hardware/software (HW/SW) partitioning remains difficult because the best partition is rarely the one with the lowest standalone execution cost; it is the one that interacts best with precedence, resource serialization, hardware area limits, and cross-partition communication. Classical exact methods quickly become impractical as graph size grows, while metaheuristics often spend substantial time exploring large binary search spaces. Recent graph-neural approaches improve scalability by turning partitioning into a differentiable graph learning problem, but most formulations still optimize placement more directly than schedule order.

This draft paper studies a repository implementation that extends differentiable graph-based HW/SW partitioning in two stages. The first stage, \texttt{diff\_gnn}, learns a relaxed hardware-assignment probability for each task node and optimizes a differentiable surrogate of makespan with area and communication penalties. The second stage, \texttt{diff\_gnn\_order}, adds explicit ordering heads, soft permutations produced through Sinkhorn normalization, and an order-aware decode procedure that couples placement and scheduling structure more tightly. The implementation also combines paper-inspired node features, learned edge weighting, differentiable local search, and a queue-based or list-scheduling-style discrete evaluator.

The goal of this manuscript is not to claim a final publishable result, but to provide a realistic six-page research-paper scaffold that mirrors the structure of the 2021 GCN-based HW/SW partitioning paper and the current 2025/2026 internal experiment draft. Using numbers already present in the repository, we show how to present both favorable and mixed evidence: on a fifteen-instance random benchmark summary the ordered model improves mean and median makespan over the baseline, while on a small eleven-node demonstration graph the benefit depends on the exact regularization setting. The resulting draft illustrates one credible paper narrative for later manual rewriting.
\end{abstract}
